\documentclass[letterpaper,11pt]{article}

\usepackage{latexsym}
\usepackage[empty]{fullpage}
\usepackage{titlesec}
\usepackage{marvosym}
\usepackage[usenames,dvipsnames]{color}
\usepackage{verbatim}
\usepackage{enumitem}
\usepackage[hidelinks]{hyperref}
\usepackage{fancyhdr}
\usepackage[brazil]{babel}
\usepackage{tabularx}
\usepackage{mathptmx}
\usepackage[utf8]{inputenc}
\usepackage[T1]{fontenc}
\input{glyphtounicode}

\pagestyle{fancy}
\fancyhf{}
\fancyfoot{}
\renewcommand{\headrulewidth}{0pt}
\renewcommand{\footrulewidth}{0pt}

% Margens
\addtolength{\oddsidemargin}{-0.5in}
\addtolength{\evensidemargin}{-0.5in}
\addtolength{\textwidth}{1in}
\addtolength{\topmargin}{-.5in}
\addtolength{\textheight}{1.0in}

\urlstyle{same}

\raggedbottom
\raggedright
\setlength{\tabcolsep}{0in}

% Formatação de seções
\titleformat{\section}{
    \vspace{-4pt}\scshape\raggedright\large
}{}{0em}{}[\color{black}\titlerule \vspace{-5pt}]

% ATS parsable
\pdfgentounicode=1

%-------------------------
% Comandos customizados
\newcommand{\resumeItem}[1]{
    \item\small{
            {#1 \vspace{-2pt}}
    }
}

\newcommand{\resumeSubheading}[4]{
    \vspace{-2pt}\item
    \begin{tabular*}{0.97\textwidth}[t]{l@{\extracolsep{\fill}}r}
    \textbf{#1} & #2 \\
    \textit{\small#3} & \textit{\small #4} \\
    \end{tabular*}\vspace{-7pt}
}

\newcommand{\resumeSubSubheading}[2]{
    \item
    \begin{tabular*}{0.97\textwidth}{l@{\extracolsep{\fill}}r}
    \textit{\small#1} & \textit{\small #2} \\
    \end{tabular*}\vspace{-7pt}
}

\newcommand{\resumeProjectHeading}[2]{
    \item
    \begin{tabular*}{0.97\textwidth}{l@{\extracolsep{\fill}}r}
    \small#1 & #2 \\
    \end{tabular*}\vspace{-7pt}
}

\newcommand{\resumeProjectEntry}[2]{
    \vspace{-2pt}\item
    \textbf{#1} \\
    \textit{\small #2}
    \vspace{-5pt}
}

\newcommand{\resumeSubItem}[1]{\resumeItem{#1}\vspace{-4pt}}

\renewcommand\labelitemii{$\vcenter{\hbox{\tiny$\bullet$}}$}

\newcommand{\resumeSubHeadingListStart}{\begin{itemize}[leftmargin=0.15in, label={}]}
\newcommand{\resumeSubHeadingListEnd}{\end{itemize}}
\newcommand{\resumeItemListStart}{\begin{itemize}}
\newcommand{\resumeItemListEnd}{\end{itemize}\vspace{-5pt}}

%-------------------------------------------
%%%%%%  CURRÍCULO COMEÇA AQUI  %%%%%%%%%%%%%%%%%%%%%%%%%%%%

\begin{document}

%----------CABEÇALHO----------
\begin{center}
\textbf{\huge \scshape João Victor Santos da Costa} \\[4pt]
\small
Uruguaiana, Rio Grande do Sul (Remoto) \quad|\quad +55 (55) 9 9217-7324 \\[2pt]
\href{mailto:victorsantos.ft18@gmail.com}{victorsantos.ft18@gmail.com} \quad|\quad
\href{https://linkedin.com/in/viquitor}{linkedin.com/in/viquitor} \quad|\quad
\href{https://github.com/Jao-Viquitor}{github.com/Jao-Viquitor}
\end{center}

%-----------RESUMO PROFISSIONAL-----------
\section{Resumo Profissional}
Engenheiro de Software com experiência em desenvolvimento \textbf{Fullstack}, especializado na criação de aplicações inteligentes utilizando \textbf{Python (FastAPI, Flask)} e frameworks modernos de frontend. Proficiente em \textbf{TypeScript, HTML5 e CSS3}, com forte base em arquiteturas de \textbf{APIs REST} e integrações complexas.
Entusiasta de Inteligência Artificial, com experiência prática no uso de \textbf{LLMs e ferramentas de Copiloto} (GitHub Copilot, Claude, Antigravity) para otimização de workflow, geração de código e arquitetura.
Sólidos conhecimentos em \textbf{SQL e bancos de dados relacionais} (PostgreSQL, MySQL), além de vivência em metodologias ágeis e boas práticas de desenvolvimento (Clean Code, SOLID). Disponibilidade imediata para atuação remota.

%-----------HABILIDADES-----------
\section{Habilidades e Tecnologias}
\begin{itemize}[leftmargin=0.15in, label={}]
\small{\item{
\textbf{Linguagens:} Python, TypeScript, JavaScript (ES6+), Java, Kotlin, SQL, HTML5, CSS3 \\
\vspace{1mm}
\textbf{Frontend:} Angular (Interesse/Aprendizado acelerado), React, Next.js, Redux, Tailwind CSS \\
\vspace{1mm}
\textbf{Backend:} FastAPI, Flask, Django, Node.js, Spring Boot \\
\vspace{1mm}
\textbf{Inteligência Artificial:} Experiência com LLMs (OpenAI, Claude), LangChain, GitHub Copilot, Antigravity Agent, Engenharia de Prompt \\
\vspace{1mm}
\textbf{Bancos de Dados:} PostgreSQL, MySQL, SQLite, MongoDB \\
\vspace{1mm}
\textbf{Ferramentas \& DevOps:} Git, Docker, Kubernetes, CI/CD (GitHub Actions), APIs RESTful \\
\vspace{1mm}
\textbf{Qualidade:} Clean Code, SOLID, Testes Unitários e de Integração, TDD \\
\vspace{1mm}
\textbf{Soft Skills:} Proatividade, Comunicação Assertiva, Protagonismo em IA, Colaboração Remota
}}
\end{itemize}

%-----------EXPERIÊNCIA-----------
\section{Experiência Profissional}
\resumeSubHeadingListStart

\resumeSubheading
{RCDNC}{Dez. 2023 -- Dez. 2025}
{Engenheiro de Software Fullstack (Foco em IA e Automação)}{Remoto}
\resumeItemListStart
    \resumeItem{\textbf{Desenvolvimento IA-Assisted:} Utilização avançada de ferramentas de IA e copilotos de código para acelerar o desenvolvimento de features, refatoração de código legado e criação de suítes de testes.}
    \resumeItem{\textbf{Backend com Python:} Implementação de microsserviços e integrações utilizando \textbf{Python (FastAPI)}, focando em performance e escalabilidade de APIs REST.}
    \resumeItem{\textbf{Integrações e APIs:} Desenvolvimento de conectores para APIs de terceiros e serviços internos, garantindo a fluidez do tráfego de dados e segurança via OAuth2/JWT.}
    \resumeItem{\textbf{Bancos de Dados:} Modelagem e otimização de queries \textbf{SQL} em bancos relacionais, assegurando a integridade e performance das aplicações.}
    \resumeItem{\textbf{Frontend Moderno:} Atuação no desenvolvimento de interfaces responsivas com \textbf{TypeScript} e frameworks SPAs, garantindo acessibilidade e experiência do usuário (UX).}
    \resumeItem{\textbf{Boas Práticas:} Mentoria de desenvolvedores e aplicação rigorosa de \textbf{SOLID} e \textbf{Clean Architecture} em projetos mobile e web.}
\resumeItemListEnd

\resumeSubheading
{ERP Vendas}{Jul. 2022 -- Out. 2022}
{Desenvolvedor Fullstack}{Remoto}
\resumeItemListStart
    \resumeItem{Manutenção e evolução de sistemas de gestão (ERP), atuando tanto no backend quanto no frontend para resolver bugs complexos e implementar novas funcionalidades.}
    \resumeItem{Otimização de consultas em bancos de dados relacionais, reduzindo o tempo de carregamento de relatórios e módulos de vendas.}
\resumeItemListEnd

\resumeSubheading
{Aroma do Sul}{Jun. 2020 -- Out. 2023}
{Gerente de Loja (Gestão e Protagonismo)}{Uruguaiana, RS}
\resumeItemListStart
    \resumeItem{Gestão de operações e marketing digital, desenvolvendo forte visão de negócio, responsabilidade e proatividade em ambientes dinâmicos.}
\resumeItemListEnd

\resumeSubHeadingListEnd

%-----------PROJETOS-----------
\section{Projetos Relevantes}
\resumeSubHeadingListStart

\resumeProjectEntry
{Virtuosa -- Processamento Inteligente de Áudio}
{Python, Librosa, FastAPI, React, PostgreSQL, Docker}
\resumeItemListStart
    \resumeItem{\textbf{Backend Inteligente:} Desenvolvimento de microsserviço em \textbf{Python} utilizando a biblioteca \textbf{Librosa} para processamento digital de sinais (DSP) e análise inteligente de BPM e padrões musicais.}
    \resumeItem{\textbf{API REST:} Exposição de dados processados via \textbf{FastAPI}, integrando o motor de análise ao ecossistema mobile/web.}
    \resumeItem{\textbf{Frontend:} Criação de dashboard administrativo com \textbf{TypeScript/React} para monitoramento de métricas e controle de usuários.}
    \resumeItem{\textbf{Infraestrutura:} Conteinerização completa do ambiente com \textbf{Docker}, facilitando o deploy e a escalabilidade dos serviços de processamento.}
    \resumeItem{Repositório: \href{https://github.com/Jao-Viquitor/virtuosa}{github.com/Jao-Viquitor/virtuosa}}
\resumeItemListEnd

\resumeSubHeadingListEnd

%-----------FORMAÇÃO-----------
\section{Formação Acadêmica}
\resumeSubHeadingListStart
\resumeSubheading
{Unopar}{Remoto}
{Bacharelado em Engenharia de Software -- Concluído em Dez. 2025}{Mar. 2024 -- Dez. 2025}
\resumeSubheading
{Universidade Federal do Pampa (Unipampa)}{Alegrete, RS}
{Bacharelado em Engenharia de Software -- Cursados semestres 1--5}{Mar. 2017 -- Set. 2022}
\resumeSubHeadingListEnd

%-----------IDIOMAS-----------
\section{Idiomas}
\begin{itemize}[leftmargin=0.15in, label={}]
    \small{\item{
    \textbf{Português:} Nativo \\
    \textbf{Inglês:} Intermediário (Leitura técnica e interação com documentação de IA) \\
    \textbf{Espanhol:} Intermediário
    }}
\end{itemize}

\end{document}
